%---------- Inleiding ---------------------------------------------------------

\section{Inleiding}%
\label{sec:inleiding}

In moderne zorglabs groeit het aantal verbonden apparaten snel. Denk dan maar aan IoT-apparatuur die continu data uitwisselen. De huidige WiFi-infrastructuren botsen echter op beperkingen zoals interferentie, beperkte dekking en uitdagingen rond schaalbaarheid en beveiliging.

5G biedt een oplossing door hogere bandbreedte en lagere latentie te hebben dan WiFi, evenals een hogere mate van mobiliteit en betere ondersteuning voor IoT-communicatie. In het zorg lab komt daarom de vraag of 5G een geschikte en veilige technologie is om bestaande en toekomstige toestellen te verbinden.

Tegelijkertijd brengt 5G ook nieuwe risico’s met zich mee. Voorbeelden zijn kwetsbaarheden in slicing-architecturen, configuratiefouten in private 5G-netwerken, mogelijke interceptie van gevoelige gezondheidsdata en de nood aan sterk device- en SIM-beheer. Voor zorginstellingen, waar dataveiligheid kritiek is, is het essentieel om deze risico’s grondig te begrijpen.

De doelgroep van dit onderzoek bestaat uit IT-beheerders, cybersecurity-specialisten en onderzoekers binnen de zorgsector die veilige 5G-connectiviteit willen realiseren.

%---------------------------------------------------------------------------
\section{Hoofdonderzoeksvraag}

\textbf{Hoe kan een private 5G-netwerkomgeving veilig en technisch haalbaar worden opgezet en geconfigureerd voor zorglabtoestellen, en welke beveiligingsmaatregelen zijn hierbij noodzakelijk?}

\vspace{1em}

%---------------------------------------------------------------------------
\section{Deelvragen -- Probleemdomein}

\begin{enumerate}
	\item Op welke manieren kunnen zorglab-toestellen technisch gekoppeld worden aan een 5G-netwerk?
	\item Welke vereisten en beperkingen hebben deze toestellen op het vlak van connectiviteit, latency, energieverbruik en protocolondersteuning?
	\item Welke beveiligingsrisico's brengt 5G-connectiviteit met zich mee binnen een zorg omgeving?
\end{enumerate}

\clearpage  % of \newpage

\begin{enumerate}
	\setcounter{enumi}{3}  % Doorgaan vanaf item 4
	\item Welke standaarden en best practices bestaan er rond beveiligde 5G- en IoT-communicatie in de zorgsector?
\end{enumerate}

\vspace{1em}

%---------------------------------------------------------------------------
\section{Deelvragen -- Oplossingsdomein}

\begin{enumerate}
	\item Welke beveiligingsmaatregelen zijn noodzakelijk om zorglab-devices veilig via private 5G te verbinden?
	\item Hoe kan toegang tot deze apparaten veilig worden geauthentiseerd  en gecontroleerd?
	\item Hoe kan het netwerkverkeer van zorglab toestellen via 5G worden gemonitord en geanalyseerd op beveiligingsrisico’s?
	\item In welke mate is private 5G technisch haalbaar (latency, stabiliteit) voor gebruik binnen het zorglab omgeving?
\end{enumerate}

%---------- Stand van zaken ---------------------------------------------------

\section{Literatuurstudie}%
\label{sec:literatuurstudie}

Onderzoek stelt verschillende risico's voor 5G-netwerken in kritieke omgevingen vast. \textcite{ENISA2022} bespreekt mogelijke kwetsbaarheden in 5G-kernarchitecturen, waaronder configuratiefouten, slicing-kwetsbaarheden en verhoogde risico's door virtualisatie. Volgens \textcite{GSMA2023} vereist 5G door zijn complexiteit extra veiligheidslagen, waaronder strengere authenticatie, integriteit van signaling en enhanced encryption. Verder tonen studies aan dat medische IoT-toestellen vaak beperkte beveiligingscapaciteiten hebben, wat risico's verhoogt bij integratie in geavanceerde netwerken zoals 5G .

De 3GPP-standaard TS~33.501 \autocite{3GPP2024} beschrijft beveiligingsarchitecturen binnen het 5G-ecosysteem, waaronder mutual authentication, SEPP-beveiliging en user-plane-encryption. \textcite{IBMSecurity2023} en andere industriële richtlijnen benadrukken het belang van device-lifecycle-management, SIM/eSIM-beveiliging, netwerksegmentatie en zero-trust-modellen voor 5G-IoT-omgevingen. Daarnaast bestaan er succesvolle implementaties van private 5G-netwerken in industriële en medische contexten, wat mogelijkheden biedt voor gecontroleerde en sterk beveiligde experimenten in een zorglab omgeving.

%---------- Methodologie ------------------------------------------------------
\needspace{15\baselineskip}

\section{Methodologie}%
\label{sec:methodologie}

\section*{Onderzoeksaanpak}

Het onderzoek wordt uitgevoerd in vijf iteratieve sprints. Door deze aanpak kunnen tussentijdse bevindingen snel worden verwerkt en kan de focus bijgestuurd worden op basis van de resultaten uit eerdere sprints indien dit nodig is. Elke sprint levert concrete resultaten op die dienen als input voor de volgende fase van het onderzoek.

\section*{Sprint 1}

In de eerste sprint wordt een uitgebreid overzicht opgesteld van verschillende 5G-technologieën en netwerkarchitecturen, samen met een lijst van mogelijke risico’s en dreigingen voor IoT-apparaten. Hiervoor worden literatuurbronnen geraadpleegd. Elke bevinding wordt systematisch gedocumenteerd in een tabel, waarin onder andere de aard van het risico, de vereiste toegangsrechten en de potentiële impact op de zorgomgeving worden weergegeven. De sprint wordt afgesloten met een risicomatrix waarin de impact en de exploitatiemoeilijkheid van de geïdentificeerde risico’s worden beoordeeld, toegespitst op de context van zorglab-toestellen.

\section*{Sprint 2}

In de tweede sprint wordt de huidige architectuur van de zorglab-toestellen en hun netwerk bestudeerd en gedocumenteerd. Dit omvat het opstellen van een overzicht van de toestellen en hun connectiviteitseisen. Daarnaast wordt gekeken of de apparaten compatibel zijn met private 5G. Ook wordt er naar connectiviteit, latency, protocolgebruik en beveiligingsvereisten gekeken. Op basis van deze analyse wordt een private 5G testomgeving opgezet en ontworpen.

\section*{Sprint 3}

In de derde sprint wordt een private 5G proof-of-concept opgezet
in een gecontroleerde testomgeving. Deze POC bestaat uit een softwarematige private 5g core, sim gebaseerde authenticatie en een zorglab-toestel of een gelijkaardig iot-apparaat. Het toestel wordt via 5G verbonden met het netwerk, waarna connectiviteit, IP-toewijzing, latency en stabiliteit worden getest.

\section*{Sprint 4}

In de vierde sprint ligt de focus op de securityanalyse en het ontwerpen van mitigaties. Aan de hand van het opgebouwde private 5G omgeving wordt een STRIDE-analyse uitgevoerd om de belangrijkste dreigingen te identificeren.

\section*{Sprint 5}

In de vijfde sprint worden alle resultaten samengebracht in een eindrapport en een presentatie.

%---------- Verwachte resultaten ----------------------------------------------

\section{Verwacht resultaat, conclusie}%
\label{sec:verwachte_resultaten}

\section*{Verwachte Resultaten}

Op basis van literatuur en initieel onderzoek wordt verwacht dat:

\begin{itemize}
	\item Private-5G technisch haalbaar is voor het verbinden van zorglab-toestellen.
	\item private 5G of netwerksegmentatie binnen 5G de meest geschikte implementatievormen blijken voor een zorglab omgeving.
	\item specifieke 5G-gerelateerde beveiligingsrisico’s zoals SIM-beheer en netwerksegmentatie relevant zullen blijken binnen een zorgcontext.
	\item zorglab-toestellen bijkomende beveiligingsmaatregelen nodig hebben.
	\item Continue monitoring en sterke authenticatie noodzakelijk zijn.
\end{itemize}
De bachelor proef resulteert daarnaast in:
\begin{itemize}
	\item een werkende private 5G proof-of-concept
	\item gedocumenteerde netwerk- en beveiligingsconfiguraties
	\item een beveiligingsanalyse van de opgebouwde testomgeving
	\item concrete en toepasbare aanbevelingen voor zorginstellingen.
\end{itemize}

\section*{Conclusie}

In deze bachelor proef wordt onderzocht of het mogelijk is om 
zorglab-toestellen op een veilige en technisch haalbare manier te 
verbinden via een private 5G-netwerkomgeving. Hiervoor wordt een 
praktische proof-of-concept opgezet in een gecontroleerde testomgeving.


Binnen deze testomgeving ligt de focus op netwerkconfiguratie en 
beveiliging, zoals SIM-gebaseerde authenticatie, netwerksegmentatie en 
monitoring. Aan de hand van deze proof-of-concept wordt de haalbaarheid 
van private 5G voor gebruik in een zorglab geëvalueerd.


De bachelor proef zal resulteren in een werkende testopstelling, 
gedocumenteerde netwerk- en beveiligingsconfiguraties en een 
beveiligingsanalyse. Deze resultaten kunnen dienen als een leidraad voor het zorglab.
\needspace{15\baselineskip}
